\chapter{Implications, proof of concept and conclusions}
\label{ch:ch8}

This final chapter translates the analytical findings of Chapters~\ref{ch:ch4}--\ref{ch:ch7} into implications for the design of elderly-oriented services and for the comparative study of ageing welfare regimes. It also introduces a proof-of-concept profiling tool based on a ten-variable subset of the SHARE battery, discusses the limitations of the analysis, and outlines directions for future work.

\section{Implications for intervention design}
\label{sec:ch8-interventions}

Each of the five Italian profiles occupies a specific position in the health / connectedness / subjective space, and each therefore suggests a different intervention logic. The mapping below is programmatic rather than prescriptive: its purpose is to make visible how the segmentation translates into actionable service design.

\emph{Fragile Resigned} (13.2\% of the Italian sample) combines physical fragility, cognitive decline and subjective acceptance of the condition. The intervention logic for this profile is \emph{outreach}: the barrier is neither economic nor informational, but motivational, since respondents do not report depressive distress and therefore are not captured by screening tools that rely on self-reported unmet need. Reach-out home visits conducted by community-health workers, with standardised isolation-screening instruments (UCLA-3, Lubben-6) administered in person, are the appropriate channel for this profile.

\emph{Fragile Depressed} (26.7\%) shares the same objective conditions as the Resigned profile but reports elevated depressive symptoms and moderate loneliness. The intervention logic is \emph{mental-health outreach with caregiver support}: the EURO-D score and the loneliness index are actionable signals that a structured intervention can address, provided it is coupled with support for informal carers (typically spouse or adult child), who in the Italian familistic regime bear the bulk of the symptomatic load.

\emph{Moderate Isolated} (26.9\%) is the hinge of the intervention portfolio. Chapter~\ref{ch:ch7} showed that this profile under-uses specialist and dentist services at sharply higher rates than matched profiles with strong social networks, but does not forgo care for cost reasons. The intervention logic is therefore \emph{navigational support and healthcare literacy}, not economic subsidy. Concrete implementations include: (i) a primary-care-led "navigator" service that proactively contacts respondents matching the profile signature and walks them through the local preventive-care calendar; (ii) a short, ten-item questionnaire administered during the annual \emph{medico di base} visit that flags the profile and activates the navigator; (iii) a community-pharmacy channel for dental-care vouchers conditioned on a simple profile screen.

\emph{Traditional Social} (7.7\%) is the profile whose social infrastructure the intervention portfolio should \emph{protect}. Policy measures here are preventive: maintenance of community centres, parish associations, neighbourhood-watch programmes, extended-family support structures. The evidence of Chapter~\ref{ch:ch7} is that the social network of this profile operates as a decentralised literacy infrastructure about the healthcare system; erosion of that infrastructure would move its members toward the Moderate Isolated profile and expose them to the same preventive-care deficit.

\emph{Connected Active} (25.4\%) is the profile best served by the current configuration. The intervention logic is \emph{preventive-service continuity}: digital-triage tools, teleconsultation channels, online booking systems, and engagement programmes (volunteer schemes, lifelong-learning courses) consolidate a profile that is already well integrated into the system.

\section{Lessons from the Swedish comparison}
\label{sec:ch8-sweden-lessons}

Chapter~\ref{ch:ch5} documented that the same six-dimensional analytical pipeline produces rhyming profile structures in Italy and Sweden, but with a systematic Swedish advantage across every positive indicator. The gap is largest at the vulnerable end of each distribution: the Italy--Sweden internet-use gap is 23 percentage points for the Connected pair but 81 percentage points for the Declining pair, and the CASP gap is 4.3 points for the Connected pair but 9.6 points for the Declining pair. This distribution-concentration pattern is diagnostic of welfare-regime architecture rather than individual endowment: the Swedish universalist system compresses the distance between the top and bottom of the ageing experience, while the Italian familistic system reproduces it.

Three lessons flow from this comparison. First, the absence of a Swedish analogue to the Italian Moderate Isolated profile suggests that universalist healthcare entry points reduce the penalty paid by individuals with low social integration. Where the healthcare system presents itself proactively --- scheduled check-ups, opt-out preventive screening, structured dental-care appointments --- the navigational-capital advantage of the Traditional Social profile is neutralised, because the system itself does what the social network does in Italy. Second, Swedish digital penetration is near-universal across the distribution and correlates with but is not reducible to socio-economic status; the policy takeaway is that digital-inclusion programmes for the elderly should be framed as universal access rather than as a targeted remediation of a minority. Third, the Swedish CASP-12 gap is structural, not composition-driven: it persists after matching on age, health and education, and therefore cannot be closed by changing who receives services, only by changing how services are organised.

\section{Proof of concept: the ten-question profiler}
\label{sec:ch8-poc}

A prototype profiling tool was developed to make the segmentation operationally usable outside the SHARE research setting. The tool is a ten-question form that asks for a small subset of the thirty-one analytical variables --- the ones whose joint answers contain the most information for classification --- and returns a predicted profile plus a visual benchmark against the Italian and Swedish distributions. The prototype, available at \texttt{v9/proof\_of\_concept/SilverItalyProfiler.jsx}, is intended as a conversation starter between the respondent and a primary-care provider, not as a diagnostic instrument: it makes the profile visible to both parties, frames the intervention logic of Section~\ref{sec:ch8-interventions}, and provides the cross-country benchmark of Chapter~\ref{ch:ch5} as a policy reference.

The ten variables were selected by iterative elimination: starting from the full 31-variable set, the variable whose removal least degraded classification accuracy was dropped at each step, until ten variables remained. The resulting subset spans all five dimensions and preserves the CASP / EURO-D / loneliness distinction that is the analytical contribution of the thesis. A deployment pilot in a municipal primary-care network would be a natural next step, beyond the scope of this thesis.

\section{Limitations}
\label{sec:ch8-limitations}

Five limitations deserve explicit mention.

First, SHARE Wave 9 is cross-sectional, so the analysis supports descriptive and associational claims rather than causal ones. Where Chapter~\ref{ch:ch7} speaks of the Moderate Isolated profile's "under-use" of discretionary healthcare services, the interpretation is that isolation and under-use co-occur in a specific configuration; the direction of causation cannot be established from W9 alone. A longitudinal follow-up across the five SHARE waves in which the relevant variables are available would be required for a causal reading.

Second, the Swedish factor solution contains a Heywood case on \emph{social\_integration} (loading 1.11; communality 1.38), as documented in Chapters~\ref{ch:ch3} and~\ref{ch:ch6}. The decision to cluster on the standardised original variables rather than on factor scores was taken partly to insulate the downstream analysis from this instability, but the loading itself remains a fragility of the Swedish FA that a larger sample or an alternative parameterisation might resolve.

Third, two Italian digital-engagement items (\emph{ac035d4}, \emph{ac035d7}) have zero variance in the Italian sample and are excluded from the Italian PCA, factor analysis and clustering. This is interpreted as an artefact of the Italian fieldwork rather than a substantive finding about non-participation, but the exclusion means that two dimensions of civic engagement are silent in the Italian segmentation.

Fourth, the healthcare utilisation analysis of Chapter~\ref{ch:ch7} relies on self-reported 12-month recall, which is subject to recall bias that is unlikely to be uniform across profiles. Respondents with stronger cognitive function (and therefore, in this sample, more connected profiles) are likely to report contacts more completely, which would attenuate --- not amplify --- the gap that Table~\ref{tab:ch7-isolati-sociali} documents. The interpretation is therefore conservative.

Fifth, the matched-pair design of Chapter~\ref{ch:ch5} maps four Italian profiles to four Swedish profiles but leaves the Italian Moderate Isolated and the Swedish Asset Rich / Wealthy Digital profiles unmatched. This is a feature, not a bug, of the segmentation --- these profiles are the welfare-regime-specific signatures that the analysis is designed to detect --- but it means that no direct cross-country $t$-test is available for those configurations.

\section{Directions for future work}
\label{sec:ch8-future}

Three research angles suggest themselves as natural follow-ups. First, replication with SHARE Wave 10 (when released) to check whether the five Italian and six Swedish profiles are stable over time, and whether the navigational-capital pattern documented in Chapter~\ref{ch:ch7} persists post-pandemic. Second, extension of the cross-country design to a Mediterranean cluster (Spain, Greece) and a Nordic cluster (Denmark, Finland) to test whether the welfare-regime gradient observed here generalises as a regime-level rather than a country-specific signature. Third, a field pilot of the ten-question profiler in a municipal primary-care network --- ideally paired with an evaluation of the navigator-service intervention proposed in Section~\ref{sec:ch8-interventions} --- to test whether the segmentation produces measurable improvements in preventive-care uptake for the Moderate Isolated profile.

\section{Conclusions}
\label{sec:ch8-summary}

The multidimensional segmentation framework proposed in this thesis identifies five empirical profiles in the Italian over-65 population and six in the Swedish one. The inclusion of the subjective block as a clustering input, rather than as a validation outcome, separates the Fragile Resigned and Fragile Depressed profiles in Italy --- a distinction that a purely objective framework collapses, as shown by the sensitivity analysis of Chapter~\ref{ch:ch6}. The cross-country comparison reveals a structural Swedish advantage on every positive indicator, largest at the vulnerable end of the distribution, and identifies three welfare-regime-specific profiles (Moderate Isolated in Italy; Asset Rich and Wealthy Digital in Sweden) that have no cross-country analogue. The healthcare-utilisation analysis establishes that the Moderate Isolated profile pays a measurable cost in terms of preventive-care access, and that the mechanism is the absence of \emph{navigational capital} provided by a dense social network, rather than the inability to afford care. The practical contribution is a profile-aware framework for the design of elderly-oriented services, instantiated in the ten-question proof of concept, that makes the social-network dimension of healthcare access operationally visible.
